\documentclass[11pt]{article}
\usepackage[margin=2.6cm]{geometry}
\usepackage{mathpazo}
\usepackage{microtype}
\usepackage{booktabs}
\usepackage{graphicx}
\usepackage{float}
\usepackage{xcolor}
\usepackage{titlesec}
\usepackage{caption}
\usepackage{listings}
\usepackage[hidelinks]{hyperref}
\emergencystretch=3em

\definecolor{accent}{HTML}{1F4E79}
\definecolor{good}{HTML}{1A7F37}
\definecolor{bad}{HTML}{B42318}
\titleformat{\section}{\Large\bfseries\color{accent}}{\thesection}{0.8em}{}
\titleformat{\subsection}{\large\bfseries\color{accent!80!black}}{\thesubsection}{0.8em}{}
\captionsetup{font=small,labelfont={bf,color=accent}}
\newcommand{\armA}{\textcolor{bad}{arm A}}
\newcommand{\armB}{\textcolor{accent}{arm B}}

\lstdefinestyle{json}{
  basicstyle=\ttfamily\footnotesize,
  backgroundcolor=\color{gray!8},
  frame=single, rulecolor=\color{gray!40}, framerule=0.4pt,
  breaklines=true, columns=fullflexible, keepspaces=true,
  showstringspaces=false, aboveskip=8pt, belowskip=8pt}
\lstset{style=json}

\newsavebox{\eliBox}
\newenvironment{eli5}
  {\par\medskip\noindent\begin{lrbox}{\eliBox}\begin{minipage}{\dimexpr\linewidth-2\fboxsep-2\fboxrule}\small}
  {\end{minipage}\end{lrbox}\noindent\fcolorbox{accent!40}{accent!4}{\usebox{\eliBox}}\par\medskip}

\title{\color{accent}\bfseries Structured DAG-Diff Mutation for CARL Reasoning Chains\\
\large The Diff Language, Its Grammar Enforcement, and the A/B Evaluation (full100)}
\author{GigaEvo}
\date{2026-07-03}

\begin{document}
\maketitle

\begin{abstract}
\noindent
We replaced free-form JSON rewriting of CARL reasoning-chain genomes with a
\emph{schema-constrained positional-slot diff}: the mutation LLM emits a typed edit whose
grammar is compiled per parent set, validated with pydantic, and applied deterministically.
In a paired 100-mutant A/B on the summarizer task --- run on the shipped grammar, in which
every wiring rule of the language is enforced by constrained decoding --- the diff arm
completed with \textbf{zero generation failures and zero structurally invalid programs
persisted}; the rewrite arm lost 5\% of its mutant budget to malformed children that each
burned a full evaluation, plus two calls to failed code extraction. Fitness is at parity.
The token bill is a trade, not a win: 11\% fewer output tokens (14\% per completed call)
against 46\% more input tokens, because the 13.7KB schema contract is charged as input on
every call. Two earlier diagnostic runs exposed the last two grammar gaps --- a
validator-checked ordering rule, then unbounded dependency arrays --- and both closures
(narrowed per-position enums; \texttt{maxItems}) are documented as design history
(\S\ref{sec:encodings}, \S\ref{sec:depfix}). This report explains the diff language in
full, then presents the A/B evidence.
\end{abstract}

\section{Motivation}
CARL chain genomes are wire-JSON documents (\texttt{ReasoningChain.to\_dict} plus platform
extras). The stock mutation path asks the LLM to rewrite the whole document as free text:
every emission re-risks the entire syntax, and a malformed child is only discovered
\emph{after} it has been persisted and has consumed a full evaluation. The hypothesis:
constraining mutation to a typed diff language --- where structural validity is enforced by
the provider's constrained decoding and a pydantic contract --- eliminates that failure class
at the cheapest possible stage.

\section{The diff language}
\label{sec:language}
This section describes the language in its shipped form --- the fixed-key \emph{object
encoding} --- which is exactly what the A/B in \S\ref{sec:experiment} ran. How the design
arrived here (two earlier diagnostic runs, each closing one grammar gap) is told in
\S\ref{sec:encodings} and \S\ref{sec:depfix}.

\subsection{ELI5}
\begin{eli5}
Imagine the child chain as a shelf of eight numbered boxes. To describe a new chain, you
fill boxes from the left. Into each box you put either a \emph{photocopy} of a step from any
parent --- naming it by its printed label, like \texttt{a2}, optionally with corrections
scribbled on top --- or a \emph{brand-new} step you write yourself. Every box also declares
which earlier boxes it reads its input from, and here is the trick: the form for box $k$
only has checkboxes for boxes $1..k-1$. Pointing at yourself or at a later box is not
``against the rules'' --- there is literally no place on the form to write it. Box~1 has no
checkboxes at all; it always reads the raw task input. When you are done, the remaining
boxes stay empty (no gaps allowed), and whatever sits in the last filled box produces the
answer that gets graded.
\end{eli5}
\noindent
The LLM never writes a chain; it fills in this form. The form is a JSON schema compiled
fresh for each mutation from the actual parents, and the provider's constrained decoding
guarantees the reply matches the form --- a violating token can never even be sampled.

\subsection{Anatomy of a diff}
A diff \emph{is} the full child chain. One payload, real shape (parents rendered as
\texttt{a1..a2} and \texttt{b1..b3}):
\begin{lstlisting}
{
  "reasoning": "Guided Innovation; insight: faithfulness - final step.
      Insert a fact-check between draft and polish: the polish step
      paraphrases away source wording that the metric scores; checking
      the draft against the source before polishing restores it.",
  "base_parent": "A",
  "slot_1": {"kind": "keep", "id": "a1"},
  "slot_2": {"kind": "new",
             "title": "Verify facts",
             "aim": "Cross-check the draft against the source text;
                     list any unsupported claim.",
             "dependencies": ["slot_1"]},
  "slot_3": {"kind": "keep", "id": "a2",
             "edits": {"aim": "Rewrite the draft into one crisp
                       sentence, keeping verified source wording."},
             "dependencies": ["slot_2"]},
  "slot_4": null, "slot_5": null, "slot_6": null,
  "slot_7": null, "slot_8": null
}
\end{lstlisting}
Field by field:
\begin{itemize}
  \item \textbf{\texttt{reasoning}} --- a falsifiable causal hypothesis for the edit, carrying
  the evidence citations (insight anchors, intra-memory clusters, donor-step ids). Prompt
  rules forbid tautologies (``verification improves faithfulness'') and diff restatements.
  \item \textbf{\texttt{base\_parent}} --- an enum of the parent letters actually present.
  This is \emph{lineage attribution}: the parent whose run metadata and task description the
  child inherits, and the lineage credit-assignment target. It does \emph{not} restrict which
  steps may be kept.
  \item \textbf{\texttt{slot\_1} \dots\ \texttt{slot\_8}} --- the child's steps, in execution
  order. \texttt{slot\_1} is required; the rest admit \texttt{null}. Filled slots must form a
  prefix (no gaps). The last filled slot is the output step --- its answer is what gets
  scored.
\end{itemize}

\subsection{The two slot forms}
Each filled slot is one of two typed objects, discriminated by \texttt{kind}:
\begin{description}
  \item[\texttt{keep}] Reuse a rendered parent step by id. \texttt{id} is an enum compiled
  from the actual parents --- parent-prefixed (\texttt{a1}, \texttt{b3}, \dots), spanning
  \emph{all} parents, so a keep may pull a donor's step directly (crossover is free).
  Optional \texttt{edits} overlays new values onto any content field
  (\texttt{title}, \texttt{aim}, \texttt{stage\_action}, \texttt{reasoning\_questions});
  unset fields keep the parent's text. \texttt{title}/\texttt{aim} edits must be non-empty
  --- the schema carries the executor's own invariant.
  \item[\texttt{new}] A fresh LLM step: required non-empty \texttt{title} and \texttt{aim};
  optional \texttt{stage\_action} and \texttt{reasoning\_questions}.
\end{description}
Both forms take the same \texttt{dependencies} field --- except in \texttt{slot\_1}, where
the field does not exist.

\subsection{Wiring: the position-narrowed dependency enum}
\texttt{dependencies} lists the earlier slots whose outputs this step consumes --- slot
positions in the \emph{new} chain, never parent ids. The schema for slot $k$ offers the enum
$\{\texttt{slot\_1},\dots,\texttt{slot\_{k-1}}\}$ and nothing else:
\begin{center}
\footnotesize
\begin{tabular}{ll}
\toprule
slot & \texttt{dependencies} may mention \\
\midrule
\texttt{slot\_1} & --- (field absent; reads the raw task input) \\
\texttt{slot\_2} & \texttt{slot\_1} \\
\texttt{slot\_3} & \texttt{slot\_1}, \texttt{slot\_2} \\
$\vdots$ & $\vdots$ \\
\texttt{slot\_8} & \texttt{slot\_1} \dots\ \texttt{slot\_7} \\
\bottomrule
\end{tabular}
\end{center}
An empty list is legal anywhere it exists --- a step may read the raw input in parallel with
slot 1. The array is also \emph{length-capped} in the grammar (\texttt{maxItems}: $k-1$ for
slot $k$): a step can never cite more distinct predecessors than exist, and a sampler that
starts looping on a dependency token hits a forced \texttt{]} after at most seven items
instead of running away. Acyclicity of the child DAG is therefore a \emph{typing} fact, not
a checked property: every edge points strictly backwards by construction.

\subsection{Every mutation move, spelled in the language}
\begin{center}
\small
\begin{tabular}{ll}
\toprule
move & spelling \\
\midrule
edit a step & \texttt{keep} it with \texttt{edits} overriding fields \\
insert a step & put a \texttt{new} slot between two \texttt{keep}s, rewire \\
delete a step & simply never fill a slot with it \\
reorder & fill the slots in the new order \\
duplicate & \texttt{keep} the same id in two slots \\
rewire only & same \texttt{keep}s, different \texttt{dependencies} \\
crossover & \texttt{keep} ids from more than one parent \\
full rewrite & every slot \texttt{new} \\
branching / parallel DAG & two slots share a dependency; a later slot joins them \\
\bottomrule
\end{tabular}
\end{center}
Two of these in miniature. Deletion is pure omission --- collapsing a 2-step parent to its
final step is a one-slot diff:
\begin{lstlisting}
{"reasoning": "...", "base_parent": "A",
 "slot_1": {"kind": "keep", "id": "a2"},
 "slot_2": null, ..., "slot_8": null}
\end{lstlisting}
Crossover is just a keep-id from another parent --- graft B's fact-extractor in front of A's
polisher:
\begin{lstlisting}
{"reasoning": "... parent:B step:b1 ...", "base_parent": "A",
 "slot_1": {"kind": "keep", "id": "b1"},
 "slot_2": {"kind": "keep", "id": "a2", "dependencies": ["slot_1"]},
 "slot_3": null, ..., "slot_8": null}
\end{lstlisting}

\subsection{What the grammar makes impossible}
The design goal: a defect class should be \emph{unrepresentable}, not detected. What an
unconstrained rewriter routinely gets wrong, this language cannot even spell:
\begin{table}[H]\centering\small
\begin{tabular}{ll}
\toprule
illegal child & why it cannot be emitted \\
\midrule
self-dependency (\texttt{slot\_k} $\to$ \texttt{slot\_k}) & slot $k$'s enum stops at \texttt{slot\_{k-1}} \\
forward dependency (\texttt{slot\_2} $\to$ \texttt{slot\_5}) & same enum narrowing \\
any dependency on slot 1 & the field does not exist there; extra keys rejected \\
dependency cycle & impossible transitively: all edges point backwards \\
keep of a nonexistent step (\texttt{a9}) & \texttt{id} enum is compiled from the actual parents \\
9-step chain & there is no \texttt{slot\_9} property \\
empty chain & \texttt{slot\_1} is required \\
empty \texttt{title}/\texttt{aim} (new or edit) & \texttt{minLength: 1} in the schema \\
runaway \texttt{dependencies} list (\texttt{slot\_1} spammed 70$\times$) & \texttt{maxItems}: $k-1$ caps the array \\
invented fields, malformed JSON & \texttt{additionalProperties: false}; constrained decoding \\
\midrule
gap fill (\texttt{slot\_3} filled, \texttt{slot\_2} null) & \emph{the} residual pydantic validator (contiguity) \\
\bottomrule
\end{tabular}
\caption{Everything above the line is enforced by the decoding grammar itself. Contiguity is
the single rule left to a validator --- and combined with the enums it also makes
``dependency on an empty slot'' impossible.}
\end{table}

\subsection{How enforcement works}
\begin{enumerate}
  \item \textbf{Per-mutation compilation.} \texttt{AllowedDagChanges.build\_schema} parses
  the actual parents and generates a pydantic model with \texttt{create\_model}: the keep-id
  enum, the \texttt{base\_parent} enum, and eight per-position slot unions are all derived
  from what really exists right now.
  \item \textbf{Constrained decoding.} The JSON schema ships to the provider as the
  structured-output contract (\texttt{method="json\_schema"}). The provider compiles it into
  a decoding grammar; tokens that would violate the schema get zero probability. The model
  cannot emit a forward reference for the same reason it cannot emit an unbalanced brace.
  \item \textbf{Portable subset.} Provider grammar compilers accept different JSON-schema
  dialects; Gemini is the strictest we probed (hard-rejects \texttt{\$ref}/\texttt{\$defs}
  and \texttt{const}). \path{gigaevo/llm/schema_compat.py} emits every schema in the
  lowest-common-denominator subset: refs inlined, \texttt{const} $\to$ single-value
  \texttt{enum}, decorative keys dropped. Nothing is provider-conditional at runtime.
  \item \textbf{Grammar budget.} Strict compilers limit construct complexity, not bytes (a
  padded 120KB schema passes; an 87KB slot-union grammar does not). Branching the slot
  models per parent blew the cliff at 2 parents $\times$ 8 slots ($\sim$27KB); the shipped
  encoding uses a \emph{single} model --- one global parent-namespaced keep-id enum,
  \texttt{base\_parent} as a plain enum field, no union discriminator --- at 13.7KB with
  $\sim$2$\times$ headroom.
  \item \textbf{Pydantic backstop.} The same model validates the reply before anything
  touches the engine; the contiguity rule lives here. A rejection costs one LLM call ---
  nothing is persisted, nothing is evaluated.
\end{enumerate}

\subsection{From diff to child}
The applier is deterministic --- the LLM's creative surface ends at the diff:
transcription walks \texttt{slot\_1..slot\_8} until the first \texttt{null}; \texttt{keep}
slots copy content fields from the referenced parent step (any parent) and overlay
\texttt{edits}; \texttt{dependencies} become sorted, deduplicated integer positions; run
metadata and the task description come from \texttt{base\_parent}; the last step is stamped
as the output step. The result is round-tripped through \texttt{ReasoningChain} as a final
assertion and emitted as canonical wire JSON --- the child is indistinguishable in format
from a hand-written genome.

\section{Implementation}

\subsection{Operator and vocabulary seam}
Two pieces, both behind existing seams:
\begin{itemize}
  \item \path{gigaevo/evolution/mutation/structured_diff.py} --- the generic engine-level
  operator \texttt{StructuredDiffMutationOperator}: it asks an \texttt{AllowedChanges}
  vocabulary to build a \texttt{DiffSchema} for the current parents, runs a
  structured-output LLM call against that schema, validates, and applies. It knows nothing
  about chains.
  \item \path{gigaevo/chains/dag_changes.py} --- \texttt{AllowedDagChanges}, the
  chain-specific vocabulary of \S\ref{sec:language}: schema compilation, parent rendering
  (the \texttt{a1}/\texttt{b2} listings the LLM sees), the applier, and the natural-language
  contract injected into the mutation prompt.
\end{itemize}

\subsection{The encoding search}
\label{sec:encodings}
Getting the diff grammar through the strictest compiler took four encodings:
\begin{table}[H]\centering\footnotesize
\begin{tabular}{llll}
\toprule
encoding & ordering rule in grammar? & size & Gemini verdict \\
\midrule
E1: tuple-union (\texttt{prefixItems}) & yes (per-slot models) & $\sim$87KB & rejected (complexity) \\
E2: cons-list recursion (\texttt{\$ref}) & yes (per-depth) & --- & rejected (\texttt{\$ref} banned) \\
E3: flat \texttt{anyOf[keep,new]} array & \textcolor{bad}{no --- pydantic validator} & $\sim$8KB & accepted --- superseded \\
E4: fixed-key object (\texttt{slot\_k}) & \textcolor{good}{yes (narrowed enums + \texttt{maxItems})} & 13.7KB & \textcolor{good}{accepted --- \textbf{shipped, ran the A/B}} \\
\bottomrule
\end{tabular}
\caption{Raw size is not the limit; construct complexity is. E4 is the encoding described in
\S\ref{sec:language} and the one the A/B ran; E3 differed only in expressing slots as an
array and checking the ordering rule post-decode (\S\ref{sec:depfix}).}
\end{table}

\subsection{Engine integration}
The engine assumed the mutator's \texttt{base\_parent} is a 1-based integer (wire-JSON
convention); diff payloads use namespace letters, so every diff mutant died pre-persist in
the first smoke (\texttt{int('A')}). Fixed test-first with \texttt{base\_parent\_index()} in
\texttt{gigaevo/evolution/engine/mutation.py} (int or namespace letter; garbage falls back to
parent 1 with a warning). Other pre-launch fixes (config-group promotion, loader pattern
clobber, prompt-fairness rewrite for arm A, token accounting) are in
\texttt{04\_issues\_log.md}.

\section{Experiment}
\label{sec:experiment}
Paired 100-mutant runs, identical seeds (4), engine config, and models: mutation
gemini-3.5-flash (OpenRouter, reasoning effort high); chain executor Qwen3-235B-Instruct
(LiteLLM proxy). Task: \texttt{problems/chains/summarizer}, ROUGE-L fitness.
\textbf{\armA{}} = stock full wire-JSON rewrite with a chain-faithful prompt;
\textbf{\armB{}} = structured slot diff on the shipped object encoding E4 of
\S\ref{sec:language}, \texttt{maxItems} included.

\section{Results}

\subsection{Validity funnel (headline)}
\begin{table}[H]\centering
\begin{tabular}{lcc}
\toprule
 & \armA{} (rewrite) & \armB{} (diff) \\
\midrule
mutation LLM calls & 107 & 106 \\
cancelled in flight at shutdown (0 tokens) & 5 & 3 \\
generation failures (pre-persist) & \textcolor{bad}{2} (no code extracted) & \textbf{\textcolor{good}{0}} \\
children persisted & 100 & 103 \\
evaluation cut by the shutdown sweep & 0 & 3 \\
\textbf{invalid after evaluation} & \textbf{\textcolor{bad}{5}} (json parse error) & \textbf{\textcolor{good}{0}} \\
valid children & 95 & 100 \\
valid yield per completed evaluation & 95.0\% & \textbf{100\%} \\
valid yield per LLM call & 88.8\% & 94.3\% \\
\bottomrule
\end{tabular}
\caption{The arms fail at different stages with very different costs. The two
``shutdown'' rows are artifacts of the fixed mutant budget, not failures: when the budget
fills, in-flight mutation calls and evaluations are swept.}
\end{table}
\noindent
\armA{} defects pass mutation silently (the operator extracts a string) and die in
evaluation: 5 mutants --- 5\% of the budget --- persisted malformed JSON (truncated or badly
quoted), each burning a mutant slot plus a full evaluation pipeline. \armB{} completed with
\textbf{zero} failures of every kind --- no schema rejections, no call errors, nothing
invalid persisted: with the wiring rules in the grammar and the array caps of
\S\ref{sec:depfix}, nothing structurally invalid can reach the population, by construction.
The 3 \armB{} children whose evaluations were shutdown-cut all parse as valid chains with
strictly-backward wiring; they are excluded from the yield denominators.

\subsection{Fitness --- parity}
Child mean 0.392 (\armA) vs.\ 0.387 (\armB); medians 0.403 vs.\ 0.400. The same 4 seeds
re-score 0.481 vs.\ 0.486 across the two runs purely from executor sampling noise, and no
same-config variance floor exists for this problem yet; \armA's single 0.551 child versus
\armB's 0.475 best child (\armB's overall best remains its own seed at 0.486) is a one-off
top sample in a single replicate per arm, so we claim parity and nothing more.
\begin{figure}[H]\centering
\includegraphics[width=\textwidth]{ab_overview.png}
\caption{Left: per-child fitness (dots) and best-so-far (line). Right: validity funnel.}
\end{figure}

\subsection{Structural exploration}
Both arms span 1--3 steps in this run: \armA{} $\{21, 51, 23\}$ across lengths 1/2/3,
\armB{} $\{15, 65, 20\}$ --- \armB{} concentrates harder on 2-step chains. The language
admits up to 8 slots and an earlier diagnostic run did produce 4- and 8-step children, but
the effect did not replicate here, so we claim no structural-exploration advantage; the
grammar's guarantee is that whatever lengths arrive are well-formed.
\begin{figure}[H]\centering
\includegraphics[width=0.7\textwidth]{ab_nsteps.png}
\caption{Chain-length distribution of valid children.}
\end{figure}

\subsection{Cost}
\begin{table}[H]\centering
\begin{tabular}{lcc}
\toprule
mutation side (gemini-3.5-flash) & \armA{} & \armB{} \\
\midrule
tokens in & 833{,}947 & 1{,}219{,}183 ($+46\%$) \\
tokens out & 804{,}569 & \textbf{714{,}864 ($-11\%$)} \\
mean output tokens per completed call & $\sim$8{,}050 & $\sim$6{,}940 ($-14\%$) \\
\bottomrule
\end{tabular}
\caption{A trade, not a win. Diffs are cheaper to \emph{emit}, but the constrained-decoding
contract is charged as input on every call: the 13.7KB schema is $\sim$3.4K tokens
$\times$ 106 calls $\approx$ 360K --- most of the input gap. Output tokens include the
model's reasoning channel (effort high), which dominates the small diff payloads and varies
run to run; earlier diagnostic \armB{} runs measured $-27\%$ and $-42\%$ output. Executor
completion tokens: 84K (\armA) vs.\ 98K (\armB).}
\end{table}
\begin{figure}[H]\centering
\includegraphics[width=\textwidth]{ab_tokens.png}
\caption{Per-mutant token spend. Left: output tokens per mutation call (dots) with a 10-call
rolling mean (line) --- \armA{} hovers at 7--9K tokens per call, \armB{} at 4--5K with a few
reasoning-channel spikes up to 49K. Zero-token points are calls that returned no completion
(\armA's extraction failures, or in-flight calls cancelled when the mutant budget was
reached). Right: cumulative mutation tokens; solid lines show the $-11\%$ output gap, dashed
lines \armB's input premium from shipping the schema with every call.}
\end{figure}

\section{Design history: two diagnostic runs, two grammar gaps closed}
\label{sec:depfix}
The grammar the A/B ran was hardened by two earlier 100-mutant diagnostic runs, each of
which exposed exactly one failure class and forced it into the schema. The requirement was
explicit throughout: fix defects in the grammar, not post-hoc.

\subsection{Gap 1: the ordering rule lived in a validator (encoding E3)}
The first run used E3 --- slots as a flat array, the \emph{dependencies-point-backwards}
rule checked by pydantic after decoding. All three of its pre-persist rejections (3/108
calls) were the same degenerate violation: slot~1 emitting
\texttt{dependencies:\ ["slot\_1"]}, a self-reference on the first slot --- precisely the
one rule the grammar did not carry. The fix is the fixed-key object encoding of
\S\ref{sec:language}: \texttt{slot\_1..slot\_8} as separate properties, each slot's
dependency enum narrowed to earlier slots, slot~1 with no dependencies field at all.
Self- and forward-references became unrepresentable; the sole surviving validator is
contiguous fill. Probes shaped the final form: Gemini rejects the per-parent-branched
variant (grammar cliff between 7 and 8 slots) and nested union discriminators, but accepts
the single-model variant --- keep-ids in one global parent-namespaced enum,
\texttt{base\_parent} a plain enum --- at 13.7KB with $\sim$2$\times$ headroom. A welcome
side effect: keeps from \emph{any} parent are now first-class, so crossover no longer
requires re-creating donor steps. Design notes:
\texttt{experiments/carl\_dag\_diff\_dependency\_rule\_fix.md}.

\subsection{Gap 2: unbounded arrays outrun the provider's enum enforcement}
The second run, on the object encoding, produced zero ordering violations --- and two
rejections of a class nobody had predicted: \emph{degenerate repetition}. The sampler
locked into a loop and emitted \texttt{"slot\_1"} $\sim$70 times inside one
\texttt{dependencies} array; deep into the list, items came back as \texttt{"slot"} and
\texttt{""} --- the provider's enum enforcement demonstrably lapses far into a long array.
The pydantic backstop caught both replies (nothing was persisted), but the calls were
wasted. The grammar fix: \texttt{maxItems}: $k-1$ on slot $k$'s dependency array
(\texttt{Field(max\_length=k-1)}). A repetition loop now hits a forced \texttt{]} after at
most seven items --- well inside the enforcement horizon --- and a step can never cite more
distinct predecessors than exist. The A/B of \S\ref{sec:experiment} ran with this cap and
recorded \textbf{zero} schema errors in 106 calls.

\subsection{Verification}
\begin{itemize}
  \item \textbf{Unit:} 32 tests in \path{tests/chains/test_dag_changes.py} --- every mutation
  archetype constructs a valid child; eleven illegal-payload classes (self-dep, forward dep,
  slot-1 dep, dangling id, gap fill, empty chain, empty aim, 9th slot, over-cap dependency
  list, the 71-item degenerate repetition payload itself, \dots) are rejected; schema-shape
  tests walk the emitted JSON schema and assert the per-position enum narrowing and the
  \texttt{maxItems} caps; a 300-case fuzz applies random schema-valid diffs and round-trips
  every child.
  \item \textbf{Suite:} targeted run green (chains + integration, 203 passed); the full
  1740-test sweep passed on the object encoding; lint clean.
  \item \textbf{Live smoke} (5 mutants, gemini-3.5-flash, same launch path as the A/B):
  5/5 mutation calls grammar-accepted, zero schema rejections, all children valid ---
  confirming Gemini accepts \texttt{maxItems} before the full run.
\end{itemize}

\section{Verdict}
The hypothesis is \textbf{supported end-to-end on the shipped grammar}: generation failures
$2\to0$, invalid persisted programs $5\to0$, valid yield per completed evaluation
$95\%\to100\%$, fitness at parity. The token bill is a trade rather than a headline win:
$-11\%$ output tokens (the pricier direction) against $+46\%$ input tokens from shipping
the 13.7KB schema with every call. Every wiring rule of the diff language is now
unrepresentable-to-break --- the A/B itself is the evidence, completing 106 calls without a
single schema error --- and the only validator left is contiguity. Next: adopt E4 as the
standing encoding for chain problems, and let longer runs test whether the 8-slot headroom
converts to structural exploration and fitness.
\end{document}
